% page 267 of the facsimile (image p-266 of the scan)
% block 162.6 x 242.0 mm ; left margin 39.4 mm
\pgc{17.780mm}{0.000mm}{
\l{\cel{54}259}
\l{}
\l{\sou{kanelo}.\cel{2}Quaza sulko longesala, mi-cilindra, qua alternas kun}
\l{\cel{3}mikra muluro ronda, o kun aristo paralela. - DEFI.}
\l{}
\l{\sou{kanetilio}. I. Filo ek oro od arjento, spulita sur longa agulo}
\l{\cel{3}de fero, ed uzata por la fabriko di la texuri brodita ye}
\l{\cel{3}oro od arjento. - II. Filo ek latuno arjentizita,qua,}
\l{\cel{3}spiral-spulite, cirkum kordo ek tripo o cirkum metala kordo,}
\l{\cel{3}donas la kordi maxim basa-sono di la violino o di la violon-}
\l{\cel{3}celo. - III. Bendo streta, ek kotonajo, di qua la bordi esas}
\l{\cel{3}garnisita ye latuna fili, uzata da la chapelifisti mulierala}
\l{\cel{3}por muntar la formo di la ornivi chapelala. - DFIRS.}
\l{}
\l{kanguruo.\cel{2}(\sou{zool.)}\cel{3}Quadripedo ek Australia, ek la klaso}
\l{\cel{3}\textquotedbl{}marsupiali\textquotedbl{},, remarkinda pro la volumino di lua kaudo sur}
\l{\cel{3}qua lu su apogas por saltar, e pro la tre granda longeso di}
\l{\cel{3}lua membri dopa. - L.\cel{1}\sou{halmaturus}. - DEFR.}
\l{}
\l{\sou{kanibalo}.\cel{2}I. La homo qua su nutras ek karno homala.- II.}
\l{\cel{3}(\sou{metaf.}) Persono feroca quale homo qua su nutras ek karno}
\l{\cel{3}homo-parta. - DEFIRS.}
\l{}
\l{kanikulo. I. La stelo di Sirius o di la Hundo, qua su levas e}
\l{\cel{3}su kushas sama-kloke kam Suno, de la 24esma di julio til la}
\l{\cel{3}24esma di agosto. - II. Periodo dum qua la vetero esas tre}
\l{\cel{3}varma, e quan produktas ta serio de dii. - EFIS.}
\l{}
\l{\sou{kanina}. (\sou{anat.}) (Dento) qua, che la homo, esas situita inter}
\l{\cel{3}la incizivi e la molari. - EFIS.}
\l{}
\l{kankro.\cel{2}(\sou{zool.)}I.Krustaceo kun kin pari de gambi, di qua la}
\l{\cel{3}gambi avana finas ye pinci, e qua marchas (egale-bone)}
\l{\cel{3}avance o des-avance. - L.\cel{1}\sou{astacus}.- II. (\sou{astron.}) Stelaro,}
\l{\cel{3}reprezentata da kankro, qua okupas la parto di zodiako, quan}
\l{\cel{3}eniras Suno ye la solstico somerala. - EF.}
\l{}
\l{kankroido.\cel{2}(\sou{patol}.) Kancero di la pelo e di la mukosi.}
\l{}
\l{kanoo.\cel{2}Kanoto pirogatra, kun extremaji pintatra, facita ek}
\l{\cel{3}ligno dina, qua movesas ordinare per pagayi. - DEFIS.}
\l{}
\l{kanono. I. (\sou{teol. katol.})l.Eklezio-lego ; decido, da la konci-}
\l{\cel{3}lo, pri la fido e pri la diciplino ; 2. (meso)-kanono :}
\l{\cel{3}parto di la mes-oficio, vorti, pregi sakramentala, quan la}
\l{\cel{3}sacerdoto dicas nelaute, de la prefaco til la \textquotedbl{}pater\textquotedbl{}. -}
\l{\cel{3}II. (\sou{Armo).}\cel{2}Cilindro ek bronzo, giso-fero, stalo, borita}
\l{\cel{3}ad-tubatre por kontenar la projektili (kuglegi, obusi, mi-}
\l{\cel{3}tralio) quin lansas la explozo da kargajo ek pulvero quan}
\l{\cel{3}onu acendas, muntita sur afusto. - DEFIRS.}
\l{}
\l{kanoniko. (\sou{eklezio katol.}) I. (\sou{olim).}\cel{1}Kleriko reguliera o seku-}
\l{\cel{3}lara, ek ordeno religiala od ek kirko katedrala o kolegiala.}
\l{\cel{3}- II. (\sou{nun)}. Kleriko sekulara qua, kun sua egali, deliberas}
\l{\cel{3}pri la aferi di la kongregaciono. - III. Kleriko sekulara qua}
\l{\cel{3}recevis de eptkopo sua titulo honorizanta. - DEFIRS.}
}
